% !TeX root = ../tesis.tex

\section{Del circuito \HaloTwo de \Mithril a \Cardano}

Era necesario cerrar la brecha entre dos entornos que, aun compartiendo la palabra \HaloTwo, imponen restricciones muy distintas. Como vimos en el capítulo anterior, el circuito \MidnightSTMCircuit de \Mithril estaba diseñado para producir un argumento sucinto sobre \BLS. Del lado de \Cardano no alcanzaba con que el argumento fuera criptográficamente correcto, sino que debía existir un camino realista para transportar el certificado, reconstruir el \transcript, ejecutar la verificación en Aiken y mantenerse por debajo de los límites simultáneos de \maxTxExUnits y \maxTxSize. Esta sección se enfoca en cómo volver la relación \STMRelation verificable \Onchain en \Cardano. Nos apoyamos en un \fork de \PlutusHaloTwoVerifierGen, en una serie de optimizaciones sobre el verificador Aiken emitido por esa herramienta y en una reorganización del flujo de transacciones del \zkBridge para lograr \StakeDistTx y \MintTx.

\subsection{El \fork de \PlutusHaloTwoVerifierGen y la adaptación a \Midnight}

La herramienta original \PlutusHaloTwoVerifierGen fue publicada por IOG como un generador de verificadores para circuitos pequeños escritos directamente sobre \texttt{halo2\_proofs}, con ejemplos como multiplicación simple, \textit{lookups} y circuitos de prueba \cite{PlutusHalo2VerifierGenRepo, IOGHalo2PlutusVerifier}. Resolvía el problema general de extraer una \textit{verification key}, modelar las expresiones relevantes del circuito y emitir un validador en Aiken usando Handlebars. Sin embargo, estaba muy lejos del caso que necesitábamos en este proyecto: el circuito \MidnightSTMCircuit no vive sobre la API estándar de \texttt{halo2\_proofs}, sino sobre el las bibliotecas de \Midnight \texttt{midnight-proofs}, \texttt{midnight-curves}, \texttt{midnight-circuits} y \texttt{midnight-zk-stdlib}. El problema no era solamente generar un validador más grande, sino extraer y reinterpretar una familia distinta de relaciones, tipos y pasos de verificación.

Por eso el \fork debió tocar casi todas las capas del pipeline, convirtiéndose en una herramienta que por un lado traduce \Midnight a Aiken, y por otro estabiliza la interfaz entre \Mithril, el \ZkBridgeOperator y la capa \Onchain de \Cardano:
\begin{itemize}[noitemsep]
	\item agregamos soporte para el circuito \MidnightSTMCircuit, con sus \textit{gadgets}, asignaciones de \textit{witness}, tipos auxiliares y \textit{runtime} específico.
	\item incorporamos un extractor nuevo para relaciones de bibliotecas de \Midnight, junto con una capa de conversión para mapear expresiones, compromisos, evaluaciones, rotaciones y pasos del protocolo a la representación intermedia que el generador ya entendía.
	\item eliminamos el backend Plinth y reorientamos toda la herramienta a Aiken.
	\item el generador pasó a exportar entidades estables consumibles por el \zkBridge. Fija qué bytes representan la prueba, qué valores se usan como \textit{public inputs}, qué parte del acumulador se preserva entre transacciones y cómo se serializa el \ReducedRedeemer que más tarde se volverá a hashear \Onchain.
\end{itemize}

Esto fue una decisión tomada al comparar los \PCS y protocolos de apertura involucrados. Tanto la herramienta \HaloTwoMultiOpen original como el circuito real de \Mithril usan compromisos \KZG \cite{Kate2010}, pero no la misma estrategia de multi-apertura. En los ejemplos clásicos del generador, se usa \GWCNineteen (que vimos en el Capítulo 2) porque es compatible con \texttt{halo2\_proofs}, mientras que las bibliotecas de \Midnight usan \HaloTwoMultiOpen. En este último caso, la agregación no es por punto, sino por \emph{conjuntos de puntos}.

Si notamos $S_j=\{z_{j,0},\ldots,z_{j,m_j-1}\}$ al $j$-ésimo conjunto, el \textit{verifier} arma una combinación lineal de compromisos
$$Q_j(X)=\sum_i x_1^i\,p_{j,i}(X) \in \poly{\Fp}_{\leq d},$$
y después necesita reconstruir el valor de esa combinación en un punto adicional $x_3$ a partir de los valores conocidos sobre todo el conjunto:
$$R_j(x_3) = \sum_{\ell = 0}^{m_j - 1}\lambda_{j,\ell}(x_3) \, Q_j(z_{j, \ell}), \qquad Z_j(X) = \prod_{\ell = 0}^{m_j - 1}(X - z_{j, \ell}),$$
donde $\lambda_{j,\ell}(x_3)$ son los coeficientes de interpolación de Lagrange. El argumento ya no entrega un witness por cada punto, sino un \textit{witness} único $\Proof$ y valores auxiliares $q_j(x_3)$. Con ellos, el \textit{verifier} reconstruye
$$f(x_3) = \sum_j x_2^j\cdot \frac{q_j(x_3)-R_j(x_3)}{Z_j(x_3)},$$
y luego forma el acumulador derecho total
$$\mathrm{final\_com}=\sum_j x_4^j Q_j(\tau)G_1+x_4^{m}f(\tau)G_1, \qquad v = \sum_j x_4^j q_j(x_3)+x_4^{m}f(x_3),$$
para terminar en el chequeo
$$\pairing{\Proof}{\tau G_2} \stackrel{?}{=} \pairing{\mathrm{final\_com} - vG_1 + x_3\Proof}{G_2}.$$

La diferencia esencial entre ambas familias es su forma algebraica. \GWCNineteen agrega aperturas \KZG ya construidas, mientras que \HaloTwoMultiOpen mueve parte del trabajo al proceso mismo de reconstruir la apertura agregada. Esta diferencia explica por qué no podíamos reutilizar sin cambios el análisis de costos de los ejemplos originales.

\subsection{Optimizaciones del validador Aiken}

Una vez obtenido un verificador correcto para \MidnightSTMCircuit, el siguiente obstáculo es económico: pasar a respetar los límites de \maxTxSize y \maxTxExUnits de \Cardano. El validador Aiken emitido por el generador ejecutaba descompresiones de puntos en $\GOne$, multiplicaciones escalares, sumas y pairings sobre \BLS. Incluso antes de considerar el flujo completo del bridge, el verificador standalone del circuito real de \Mithril quedaba fuera del presupuesto de CPU.

Las primeras mejoras se hicieron sobre el código Aiken generado, no sobre el protocolo del \zkBridge. En el caso del \MSM izquierdo se reemplazó la evaluación genérica basada en \texttt{foldl} y listas de \texttt{MSMElement} por una cadena \textit{inline} de llamadas a \texttt{addG1} y \texttt{scaleG1}, eliminando \textit{overhead} de la máquina CEK de Aiken, y también omitimos multiplicaciones escalares redundantes cuando el coeficiente es la identidad. La misma idea se extendió luego al \MSM derecho.

Sobre ese mismo \MSM derecho, eliminamos dos fuentes de costo adicionales. Por un lado, el \textit{roundtrip} \texttt{compress/decompress} del punto \texttt{vanishing\_g}: una vez computado como elemento de $\GOne$, no había motivo para reserializarlo como \texttt{ByteArray} sólo para volver a descomprimirlo dentro del evaluador genérico del \MSM. Por otro, la descompresión duplicada de los puntos $W_i$, que aparecían tanto en el lado izquierdo como en el derecho. Deduplicar esas descompresiones no altera la semántica de la verificación, pero sí evita repetir llamadas costosas a \texttt{bls12\_381\_G1\_uncompress}.

También eliminamos términos algebraicamente muertos del cálculo del acumulador del polinomio \textit{vanishing}. En particular, el patrón \texttt{scaleG1(zero, xn)} seguido de una suma con el primer compromiso del bloque es innecesario, porque para todo $s \in \Fr$ valen $s \cdot \OO = \OO$, y $\OO + P = P$. Junto con pequeñas mejoras adicionales, esto produjo ahorros medibles de CPU en todos los circuitos de referencia y permitió que ejemplos pesados del repositorio original, como \texttt{atms\_with\_lookups}, entraran en presupuesto.

Sin embargo, ese conjunto de optimizaciones no resolvió por sí mismo el caso \MidnightSTMCircuit. El verificador de \Mithril mejoró, pero siguió siendo demasiado costoso para una sola transacción. Por eso la solución final no fue simplemente ``emitir mejor código'', sino ver cómo reorganizar el flujo para que el costo pesado de verificar un certificado de \Mithril se pague una sola vez y luego pueda ser consumido por la cadena de \StakeDistribution y por \MintTx sin reejecutar el mismo chequeo.

Aún si el presupuesto de CPU fuera suficiente, cargar en cada transacción los scripts grandes del verificador \STM, del \TxsUpdaterMintingValidator y del \MintingValidator volvía inviable el flujo completo al superar el \maxTxSize de \Cardano. Para esto, aprovechamos los \ReferenceScripts introducidos en el \CIP-33 \cite{CIP0033} para persistir scripts pesados \Onchain y referenciarlos desde las transacciones posteriores en lugar de reenviarlos como parte del \textit{witness set}. Podría ser sólo un detalle de implementación, pero fue estrictamente necesario para que el \zkBridge final entre efectivamente en el modelo \eUTxO de \Cardano. En el flujo concretamente implementado, esa decisión aparece en tres transacciones:
\begin{itemize}[noitemsep]
    \item \PublishPhaseOneReferenceScriptTx,
    \item \PublishMintingTxsUpdaterSpendReferenceScriptTx,
    \item \PublishBridgeMintingReferenceScriptTx.
\end{itemize}

Esta decisión de partir el flujo y publicar scripts pesados responde al conjunto de builtins disponible hoy. Si prospera \CIP-133, que propone una \textit{builtin} nativa de \MSM sobre \BLS \cite{CIP0133, PlutusMSMBenchRepo}, una parte importante del bloque lineal que hoy domina el costo del verificador podría colapsar en una sola operación, reduciendo drásticamente tanto el costo algebraico y el \textit{overhead} CEK, y permitiría hacer la verificación de un certificado de \Mithril sin una transacción de dos fases.

\subsection{Partición de la verificación en dos transacciones}

Como seguíamos fuera del presupuesto de \maxTxExUnits, debimos partir la validación de un argumento \STM de un certificado de \Mithril en dos transacciones consecutivas. En \MidnightSTMCircuit, el trabajo costoso se concentra en el acumulador derecho (recordemos que estamos en \HaloTwoMultiOpen), es decir, en la combinación de compromisos por \textit{point sets}, potencias de $x_1$ y $x_4$, reconstrucciones interpoladas y el término final que incorpora $f(x_3)$. Ese acumulador es el bloque lineal dominante de \ExUnits.

Hicimos el siguiente modelo algebraico usado para diseñar el \textit{split}. Para cada conjunto de puntos $S_j$, el \textit{verifier} arma una combinación lineal $Q_j$ de compromisos usando potencias de $x_1$, evalúa la parte interpolada $R_j(x_3)$, y luego combina los distintos conjuntos con potencias de $x_4$ para formar el acumulador derecho. Si denotamos por $I_j$ al conjunto de índices de compromisos de $S_j$, entonces partir el verificador equivale a escribir $I_j = I_j^{(1)} \sqcup I_j^{(2)}$ y descomponer (por linealidad) tanto $Q_j$ como $R_j(x_3)$ y el acumulador final en una suma de una parte de fase 1 y una parte de fase 2. La primera transacción no ``verifica media prueba'' en un sentido informal, sino que computa un prefijo algebraicamente bien definido del mismo acumulador que la segunda transacción terminará de cerrar.

Para este diseño, debemos decidir cómo transportar un estado intermedio compacto. En lugar de persistir la prueba completa o de recomputar todo en la segunda transacción, la fase 1 produce un \textit{datum} enriquecido que preserva exactamente lo necesario para retomar el cálculo: un acumulador parcial comprimido, tres desafíos del \transcript y un escalar adicional. En la implementación, esa estructura aparece como \PhaseOneState, mientras que la segunda transacción recibe un \ReducedRedeemer que contiene sólo las contribuciones restantes necesarias para completar el bloque lineal y ejecutar el chequeo final de pairing.

El punto de corte finalmente se eligió , sino por medición directa. Aunque el análisis preliminar sugería un corte algo más tardío, al instrumentar por separado \PhaseOneSetupTx y \PhaseTwoVerifyTx vimos que la fase 1 absorbía más trabajo fijo del esperado y que el mejor equilibrio real aparecía antes en la entrada 17 del primer conjunto de puntos. Como los cuatro conjuntos de puntos de \MidnightSTMCircuit tienen tamaños $44, 6, 3, 2$, la implementación final deja a la fase 1 las contribuciones de los primeros 17 puntos del primer conjunto, y a la fase 2 todo lo restante. Esta partición fue la que produjo el mejor balance medido entre ambas transacciones (en poco más de 7 billones de unidades), y lo hizo viable.

\subsection{Validez del transcript y \textit{binding} entre fases}
%TODO: Mejorar esta subseccion
Partir la verificación en dos transacciones introduce una obligación adicional: garantizar que la fase 2 continúa exactamente el mismo \transcript y el mismo objeto de prueba que comenzó a procesar la fase 1. Si esa continuidad fallara, aparecería una superficie de ataque obvia: usar una fase 1 barata para fijar un estado y luego completar la fase 2 con un \redeemer distinto o con un conjunto incompatible de compromisos. Toda la lógica del datum intermedio apunta precisamente a cerrar esa posibilidad.

En la primera transacción, \PhaseOneSetupTx recibe la prueba serializada y las dos entradas públicas del circuito, reconstruye el prefijo barato del verificador y deja persistido un datum de tipo \PhaseOneState dentro del \UTxO que luego consumirá la fase 2. Ese estado contiene \texttt{rhs\_prefix}, que es un acumulador parcial comprimido en $\GOne$; \texttt{reduced\_hash}, que fija el hash esperado del \ReducedRedeemer; los desafíos \texttt{x1}, \texttt{x3} y \texttt{x4}; y el escalar \texttt{v}, ya precomputado. Además, el datum extendido del \UTxO de fase 2 conserva un \texttt{statement\_hash}, que en la familia de circuitos \Lagrange coincide directamente con la segunda entrada pública del circuito, es decir, con el mensaje firmado por el certificado de \Mithril.

Es importante notar qué no se persiste. El desafío $x_2$ no reaparece en el datum porque su información ya quedó absorbida en los cálculos que produce la fase 1: una vez reconstruidos $f(x_3)$ y $v$, la segunda transacción ya no necesita rederivar $x_2$ para cerrar el chequeo. Esta decisión reduce el tamaño del estado intermedio y, al mismo tiempo, hace explícita una propiedad conceptual relevante: el datum no guarda ``todo el transcript'', sino solamente la porción de estado necesaria para volver determinista la continuación de la verificación.

El segundo mecanismo de \textit{binding} es el hash del \ReducedRedeemer. La fase 1 serializa canónicamente ese \redeemer reducido, computa \BLAKETwoSTwoFiveSix y persiste el resultado como \texttt{reduced\_hash}. Ese mismo valor se usa además como nombre del NFT emitido en fase 1 y también del NFT que la fase 2 mintea cuando la verificación final tuvo éxito. La consecuencia es fuerte: el estado intermedio, el activo que identifica el \UTxO de fase 2 y el recibo final de prueba verificada quedan todos ligados al mismo objeto serializado. Si la fase 2 intentara presentar un \ReducedRedeemer distinto, el rehash fallaría antes de ejecutar la parte cara del verificador.

El \textit{binding} final se da entre el \texttt{statement\_hash} y el certificado de \Mithril. La fase 2, cuando termina con éxito, no sólo consume el \UTxO intermedio sino que crea un \ProofReceipt con un datum \texttt{ProofReceiptDatum \{ statement\_hash \}}. Ese valor es el mismo que más tarde leerán los validadores de \StakeDistribution y de \MintTx para comprobar que el certificado presentado lleva exactamente el \textit{signed\_message} que ya fue autenticado por \MidnightSTMCircuit. Además, ese \texttt{statement\_hash} no queda respaldado meramente por estar escrito en un datum: queda respaldado por el hecho de que el \UTxO fue creado junto con el NFT minteado bajo la \texttt{PolicyId} de la fase 2, y los validadores downstream buscan precisamente \ProofReceipt marcados por esa política. Así, el \ProofReceipt no es sólo un portador de datos, sino una evidencia autenticada por los scripts que ejecutaron la segunda mitad del verificador.

\subsection{Conexión de \UTxOs con \StakeDistTx y \MintTx}

El \ProofReceipt anteriormente generado es la pieza que conecta el verificador STM con el resto del bridge. En el subflujo de \StakeDistribution, la transacción \StakeDistributionGenesisTxTemplate consume el \ProofReceipt correspondiente al dominio \texttt{stake\_distribution\_genesis} y usa su \texttt{statement\_hash} para anclar el certificado génesis que creará el primer \StakeDistributionUTxO. Luego, \StakeDistributionStandardTxTemplate consume el recibo del dominio \texttt{stake\_distribution\_standard} junto con el \StakeDistributionUTxO padre, verifica la continuidad de la cadena de certificados y emite el nuevo \StakeDistributionUTxO canónico que el resto del protocolo considerará vigente. De esta manera, el estado persistente del bridge en \ChainB depende de una cadena de \UTxOs que sólo puede avanzar si cada certificado nuevo está ligado a un \ProofReceipt de verificación \STM ya aceptado \Onchain.

En el camino $\LockTx \rightarrow \MintTx$, la transacción \BridgeMintTxTemplate consume el \ProofReceipt del dominio \texttt{cardano\_transactions}, consume también el \TxsUpdaterUTxO y referencia el \StakeDistributionUTxO vigente (como \ReferenceInput, ver \CIP-31 \cite{CIP0031}). La separación de roles es aquí muy precisa. El \ProofReceipt de \texttt{cardano\_transactions} prueba que un \TransactionSnapshotCertificate ha pasado la verificación \STM y fija el \texttt{statement\_hash} que debe coincidir con el \textit{signed\_message} del certificado. El \StakeDistributionUTxO aporta el certificado padre confiable contra el cual debe encadenarse ese \TransactionSnapshotCertificate. Y el \TxsUpdaterMintingValidator garantiza que la transacción bloqueante de \ChainA no pueda reutilizarse para mintear dos veces, reemplazando la nueva raíz de \SMT en el $\TxsUpdaterUTxO^*$ actualizado.

Recién sobre esa base, el argumento \Groth de pertenencia de la \LockTx dentro del \TransactionSnapshot certificado adquiere el significado que el \zkBridge necesita: argumenta inclusión en un snapshot autenticado por un certificado de \Mithril ya enlazado a la \CertificateChain mantenida \Onchain.

